% =============================================================================
%  Notas del profesor — Sesión 09: Examen Parcial 01
%  Programación Avanzada — Universidad Panamericana
%  Compilar: pdflatex sesion09_examen_profesor.tex
% =============================================================================
\documentclass[12pt, oneside]{article}
\usepackage[profesor]{progavanzada}

\begin{document}

\portada{Sesión 09}{Examen Parcial 01}{2026}

\tableofcontents
\newpage

% =============================================================================
\section{Resumen de la sesión}
% =============================================================================

Esta sesión es el primer examen parcial.  Cubre sesiones 01–08.
El examen se resuelve \textbf{a mano, en papel, sin computadora}.
No incluye escribir ni interpretar código fuente; el enfoque es
\textbf{teórico y matemático}.

\begin{notaprofesor}[Logística]
  \begin{itemize}
    \item \textbf{Duración:} 80 minutos (duración estándar de una sesión).
    \item \textbf{Material permitido:} pluma o lápiz.  Sin calculadora.
    \item \textbf{No permitido:} celular, calculadora, apuntes.
    \item \textbf{Entregar:} hojas con nombre completo, matrícula y fecha.
  \end{itemize}
\end{notaprofesor}

% =============================================================================
\section{Distribución temática recomendada}
% =============================================================================

El examen está diseñado para 40 puntos totales.
La distribución refleja el tiempo de clase dedicado a cada unidad y la
dificultad conceptual de cada tema.

\begin{center}
\begin{tabular}{clcc}
  \toprule
  \textbf{Sesión(es)} & \textbf{Tema} & \textbf{Puntos} & \textbf{\%} \\
  \midrule
  01–02 & Computadora, C++ y jueces en línea (conceptual) & 4  & 10\% \\
  03    & Sistemas de numeración y conversiones            & 10 & 25\% \\
  04    & Complemento a 2, ASCII, fracciones e IEEE 754   & 12 & 30\% \\
  05    & Operadores de bits: tablas de verdad y evaluación & 8 & 20\% \\
  06–08 & Técnicas de bit manipulation (conceptual y aplicada) & 6 & 15\% \\
  \bottomrule
  & \textbf{Total} & \textbf{40} & \textbf{100\%} \\
  \bottomrule
\end{tabular}
\end{center}

\begin{notaprofesor}[Criterio de diseño]
  Las conversiones y el complemento a 2 son los temas con mayor peso porque:
  (1) son la base de todo lo demás, (2) se pueden evaluar con preguntas
  graduadas en dificultad, y (3) el proceso a mano es comprobable paso a paso.
  Los temas conceptuales (sesiones 01–02) tienen peso menor porque no tienen
  componente matemática rica, pero sirven para diferenciar a quienes leyeron
  y quienes no.
\end{notaprofesor}

\newpage

% =============================================================================
\section{Tipos de preguntas recomendadas}
% =============================================================================

Dado que el examen es \textbf{a mano, sin código}, los formatos más efectivos
son los siguientes:

% -----------------------------------------------------------------------------
\subsection{Conversión de bases (10–25 puntos sugeridos)}
% -----------------------------------------------------------------------------

Son las preguntas más objetivas y fáciles de calificar.
Recomendación: pide \textbf{al menos una conversión cruzada} (e.g., octal → hex)
para verificar que el alumno domina el binario como lenguaje puente, no solo
las conversiones directas.

\begin{notaprofesor}[Tipos de conversión recomendados]
  \begin{itemize}
    \item Decimal $\to$ binario (divisiones sucesivas).
    \item Binario $\to$ decimal (suma de pesos).
    \item Hexadecimal $\to$ binario (grupos de 4 bits) — muy rápida, evalúa si
          memorizaron la tabla.
    \item Binario $\to$ hexadecimal (grupos de 4 bits desde la derecha).
    \item Octal $\to$ hexadecimal \textbf{vía binario} — evalúa comprensión del método.
    \item Base arbitraria (base 4, base 5) $\to$ decimal — evalúa comprensión general
          del sistema posicional.
  \end{itemize}
  Evita pedir decimal $\to$ octal directo (tedioso y sin valor pedagógico extra).
\end{notaprofesor}

% -----------------------------------------------------------------------------
\subsection{Complemento a 2 (8–12 puntos sugeridos)}
% -----------------------------------------------------------------------------

\begin{notaprofesor}[Preguntas de alto valor diagnóstico]
  \begin{enumerate}
    \item \textbf{Representar un negativo en $n$ bits} (mostrar el proceso).
          Detecta si el alumno aplica el algoritmo o solo memoriza el resultado.
    \item \textbf{Interpretar un patrón de bits con signo}.
          El bit más significativo es 1 → el alumno debe negar el patrón para
          encontrar el valor.  Error común: tratar el patrón como entero sin signo.
    \item \textbf{Calcular el rango} de un tipo de $n$ bits con y sin signo.
    \item \textbf{Pregunta trampa}: ``¿Cuánto vale \texttt{\textasciitilde(-128)}
          en 8 bits?''  Respuesta: $127$.  Verifica si entienden que $\texttt{\textasciitilde x} = -(x+1)$.
  \end{enumerate}
\end{notaprofesor}

% -----------------------------------------------------------------------------
\subsection{ASCII (4–6 puntos sugeridos)}
% -----------------------------------------------------------------------------

\begin{notaprofesor}[Formatos útiles]
  \begin{itemize}
    \item \textbf{Decodificar un mensaje en hex}: secuencia de 4–8 bytes.
          Elegir palabras en mayúsculas o con dígitos para que el alumno use
          la tabla, no solo memorice.
    \item \textbf{Conversión mayúscula $\leftrightarrow$ minúscula sin tabla}:
          pedir que expliquen qué bit cambia y cuál es la operación de bits
          equivalente (\texttt{c | 0x20} para pasar a minúscula,
          \texttt{c \& 0xDF} para pasar a mayúscula).
    \item \textbf{Carácter dígito $\to$ entero}: \texttt{'7' - '0' = 7} o
          \texttt{'7' \& 0x0F = 7}.  Evalúa si entienden la estructura del
          nibble en los dígitos ASCII.
  \end{itemize}
\end{notaprofesor}

% -----------------------------------------------------------------------------
\subsection{Fracciones binarias e IEEE 754 (6–8 puntos sugeridos)}
% -----------------------------------------------------------------------------

\begin{notaprofesor}[Preguntas de diferenciación]
  \begin{enumerate}
    \item \textbf{Convertir una fracción a binario} mostrando la tabla de
          multiplicaciones por 2.  Elegir una exacta y una periódica.
          Fracciones útiles: $3/8$ (exacta), $1/6$ (periódica), $5/16$ (exacta),
          $1/3$ (periódica clásica).
    \item \textbf{Sin hacer la conversión, determinar si una fracción es exacta en
          binario} por análisis de factores primos del denominador.
          Evalúa comprensión profunda vs.\ memorización.
    \item \textbf{Analizar un patrón IEEE 754}: dado el patrón de 32 bits,
          identificar signo, calcular el exponente real ($e - 127$) y
          reconstruir el valor.  No hay que conocer todos los detalles del
          estándar — basta entender la estructura.
    \item \textbf{Pregunta conceptual}: ``¿Por qué \texttt{0.1 + 0.2 $\neq$ 0.3}?''
          Una respuesta correcta menciona: (a) $0.1_{10}$ es periódico en binario,
          (b) se trunca al almacenarlo, (c) los errores se acumulan.
  \end{enumerate}
\end{notaprofesor}

% -----------------------------------------------------------------------------
\subsection{Operadores de bits (8–10 puntos sugeridos)}
% -----------------------------------------------------------------------------

\begin{notaprofesor}[Diseño de preguntas]
  \begin{itemize}
    \item \textbf{Evaluar expresión paso a paso} mostrando los 8 bits en cada
          etapa.  Incluir al menos: un desplazamiento, un NOT, y un AND o XOR.
          Ejemplo: \texttt{(\textasciitilde(n >> k)) \& m}.
    \item \textbf{Completar tabla}: dar \texttt{a} y \texttt{b} en binario,
          pedir \texttt{a \& b}, \texttt{a | b}, \texttt{a \textasciicircum{} b}.
    \item \textbf{Identificar qué hace una expresión}: mostrar la expresión y
          pedir describir en palabras qué hace.  Ejemplos:
          \texttt{n \& (n-1)}, \texttt{n \& -n}, \texttt{x | (1 << k)}.
    \item \textbf{Verificar la precedencia}: dar \texttt{a \& b | c} vs.\
          \texttt{a \& (b | c)} y pedir calcular ambos con valores concretos.
          Evidencia si los alumnos usan paréntesis o confían en la precedencia
          de memoria.
  \end{itemize}
\end{notaprofesor}

% -----------------------------------------------------------------------------
\subsection{Técnicas de bit manipulation (4–6 puntos sugeridos)}
% -----------------------------------------------------------------------------

\begin{notaprofesor}[Enfoque conceptual — sin código]
  Dado que el examen no incluye código, las técnicas se evalúan de dos formas:
  \begin{enumerate}
    \item \textbf{Aplicar la técnica a un número concreto} y mostrar el resultado
          con representación binaria.  Ejemplo: dado \texttt{x = 0b10110100},
          apagar el bit 5 y mostrar el resultado.
    \item \textbf{Explicar con palabras} qué hace la expresión y por qué.
          Ejemplo: ``¿Qué hace \texttt{n \& (n-1)} a los bits de \texttt{n}?
          Demuéstralo con $n = 24$.''
  \end{enumerate}
  Los problemas de la sesión 07 (Chasquidos Mágicos, separación de bits)
  pueden usarse como preguntas de razonamiento: ``dado $n=4$ y $c=47$,
  ¿cuál es el estado de la lámpara? Justifica sin calcular $47$ en binario
  desde cero --- razona sobre el bit $n-1$''.
\end{notaprofesor}

\newpage

% =============================================================================
\section{Preguntas de ejemplo con respuestas}
% =============================================================================

Las siguientes preguntas son representativas de lo que puede ir en el examen.
Se incluyen respuestas completas y criterios de calificación parcial.

% -----------------------------------------------------------------------------
\subsection{Conversiones}
% -----------------------------------------------------------------------------

\textbf{P1.} Convierte $219_{10}$ a binario.

\begin{notaprofesor}[Respuesta P1 — 2 pts]
  Divisiones sucesivas:
  \begin{center}
  \begin{tabular}{r@{$\;\div 2 =\;$}l@{\quad}l}
    219 & 109 & resto \textbf{1} \\
    109 & 54  & resto \textbf{1} \\
    54  & 27  & resto \textbf{0} \\
    27  & 13  & resto \textbf{1} \\
    13  & 6   & resto \textbf{1} \\
    6   & 3   & resto \textbf{0} \\
    3   & 1   & resto \textbf{1} \\
    1   & 0   & resto \textbf{1} \\
  \end{tabular}
  \end{center}
  Resultado: $219_{10} = 1101\,1011_2$.
  Verificación: $128+64+16+8+2+1 = 219$. \checkmark

  \textbf{Calificación parcial:} 1 pt si el proceso es correcto pero hay
  un error aritmético en una división.  0 pts si solo escribe la respuesta
  sin proceso.
\end{notaprofesor}

\medskip
\textbf{P2.} Convierte $\mathtt{6C3}_{16}$ a octal.

\begin{notaprofesor}[Respuesta P2 — 3 pts (1 por etapa)]
  Paso 1 — Hex $\to$ Binario:
  \[
    6 = 0110, \quad C = 1100, \quad 3 = 0011
    \;\;\Rightarrow\;\; 0110\,1100\,0011_2
  \]
  Paso 2 — Binario $\to$ Octal (grupos de 3 desde la derecha):
  \[
    \underbrace{011}_{3}\;
    \underbrace{011}_{3}\;
    \underbrace{000}_{0}\;
    \underbrace{011}_{3}
    \;\;\Rightarrow\;\; 3303_8
  \]
  Verificación: $3\cdot512 + 3\cdot64 + 0\cdot8 + 3 = 1536+192+3 = 1731$.
  $6\cdot256 + 12\cdot16 + 3 = 1536+192+3 = 1731$.\ \checkmark

  \textbf{Error frecuente:} intentar dividir \texttt{0x6C3} entre 8
  directamente — tedioso y propenso a errores.  Si el alumno llega al
  resultado correcto por ese camino, vale los puntos completos.
\end{notaprofesor}

% -----------------------------------------------------------------------------
\subsection{Complemento a 2}
% -----------------------------------------------------------------------------

\textbf{P3.} Representa $-53$ en \textbf{8 bits complemento a 2}.
Muestra el proceso.

\begin{notaprofesor}[Respuesta P3 — 3 pts]
  \begin{center}
  \begin{tabular}{ll}
    $53_{10}$ en binario  & \texttt{0011 0101} \\
    Invertir bits          & \texttt{1100 1010} \\
    Sumar 1                & \texttt{1100 1011} \\
  \end{tabular}
  \end{center}
  $-53_{10} = \texttt{1100\,1011}_2$.
  Verificación: $2^8 - 53 = 203 = 128+64+8+2+1 = \texttt{1100\,1011}$. \checkmark

  \textbf{Calificación:} 1 pt proceso (inversión correcta), 1 pt suma 1,
  1 pt resultado final.  Error en la conversión de 53 a binario descuenta 1 pt
  del primer paso pero puede haber crédito en los siguientes si el proceso
  es correcto.
\end{notaprofesor}

\medskip
\textbf{P4.} ¿Qué valor decimal tiene \texttt{1010\,0011} interpretado
como entero con signo de 8 bits?

\begin{notaprofesor}[Respuesta P4 — 3 pts]
  El bit más significativo es 1: es \textbf{negativo}.
  Aplicar el algoritmo para obtener el valor absoluto:
  \begin{center}
  \begin{tabular}{ll}
    \texttt{1010 0011} & (valor dado) \\
    Invertir           & \texttt{0101 1100} \\
    Sumar 1            & \texttt{0101 1101} = $64+16+8+4+1 = 93$ \\
  \end{tabular}
  \end{center}
  El valor es $-93$.

  \textbf{Calificación:} 1 pt por identificar que es negativo, 1 pt por
  aplicar el proceso, 1 pt por el resultado.

  \textbf{Error frecuente:} calcular $128+32+2+1 = 163$ y reportar $+163$.
  Esto indica que el alumno no sabe distinguir con/sin signo.  0 pts.
\end{notaprofesor}

% -----------------------------------------------------------------------------
\subsection{ASCII}
% -----------------------------------------------------------------------------

\textbf{P5.} Decodifica: \texttt{44 45 43 49 4D 41 4C}

\begin{notaprofesor}[Respuesta P5 — 2 pts (0.5 por cada carácter correcto con proceso)]
  Usando la tabla ASCII (rango mayúsculas 0x41–0x5A):
  \begin{center}
  \begin{tabular}{ccc}
    \toprule
    Hex & Dec & Carácter \\
    \midrule
    44 & 68 & D \\
    45 & 69 & E \\
    43 & 67 & C \\
    49 & 73 & I \\
    4D & 77 & M \\
    41 & 65 & A \\
    4C & 76 & L \\
    \bottomrule
  \end{tabular}
  \end{center}
  Resultado: \textbf{DECIMAL}.
\end{notaprofesor}

\medskip
\textbf{P6.} \texttt{0x6D} es una letra minúscula.  Sin consultar la tabla
ASCII, determina la mayúscula y su código hexadecimal.  Explica con
operaciones de bits.

\begin{notaprofesor}[Respuesta P6 — 3 pts]
  La diferencia entre mayúscula y minúscula es el \textbf{bit 5} (valor 32 = 0x20).
  Las minúsculas tienen bit 5 encendido (\texttt{011xxxxx}); las mayúsculas,
  apagado (\texttt{010xxxxx}).

  Para pasar de minúscula a mayúscula: \textbf{apagar el bit 5}:
  \[
    \texttt{0x6D \& \textasciitilde 0x20}
    = \texttt{0x6D \& 0xDF}
    = \texttt{0110\,1101} \;\&\; \texttt{1101\,1111}
    = \texttt{0100\,1101}
    = \texttt{0x4D}
  \]
  \texttt{0x4D} corresponde a la letra \textbf{M}.
  Verificación: \texttt{0x6D} es \texttt{'m'} y \texttt{0x4D} es \texttt{'M'}.
  \checkmark

  \textbf{Calificación:} 1 pt por identificar el bit 5, 1 pt por la operación,
  1 pt por el resultado correcto.
\end{notaprofesor}

% -----------------------------------------------------------------------------
\subsection{Fracciones binarias}
% -----------------------------------------------------------------------------

\textbf{P7.} Convierte $\dfrac{1}{6}$ a binario.  ¿Es exacta o periódica?

\begin{notaprofesor}[Respuesta P7 — 3 pts]
  $1/6 \approx 0.1\overline{6}$.  Multiplicaciones por 2:
  \begin{center}
  \begin{tabular}{r@{$\;\times 2 =\;$}ll}
    $0.1\overline{6}$ & $0.\overline{3}$  & bit \textbf{0} \\
    $0.\overline{3}$  & $0.\overline{6}$  & bit \textbf{0} \\
    $0.\overline{6}$  & $1.\overline{3}$  & bit \textbf{1}, resto $0.\overline{3}$ \\
    $0.\overline{3}$  & $0.\overline{6}$  & bit \textbf{0} \quad (estado repetido) \\
    $\vdots$          & $\vdots$          & $\vdots$ \\
  \end{tabular}
  \end{center}
  El patrón que se repite es \texttt{01} (a partir del tercer bit):
  $\dfrac{1}{6} = 0.00\overline{10}_2$.

  Es \textbf{periódica}.  La razón: $6 = 2 \times 3$ y el factor 3 no es
  potencia de 2, por lo tanto no tiene representación exacta en base 2.

  \textbf{Calificación:} 1 pt tabla de multiplicaciones, 1 pt identificar patrón
  periódico, 1 pt justificación por factores primos.
\end{notaprofesor}

\medskip
\textbf{P8.} Sin hacer la conversión, determina si $\dfrac{7}{16}$ y
$\dfrac{3}{14}$ son exactas en binario.  Justifica.

\begin{notaprofesor}[Respuesta P8 — 2 pts]
  \begin{itemize}
    \item $\dfrac{7}{16}$: denominador $16 = 2^4$. Solo factor primo 2.
          \textbf{Exacta en binario.}
    \item $\dfrac{3}{14}$: denominador $14 = 2 \times 7$.
          El factor 7 no es potencia de 2.
          \textbf{No exacta — periódica.}
  \end{itemize}
\end{notaprofesor}

% -----------------------------------------------------------------------------
\subsection{Operadores de bits}
% -----------------------------------------------------------------------------

\textbf{P9.} Evalúa \texttt{(\textasciitilde(30 >> 3)) \& 19}.
Muestra los 8 bits en cada paso.

\begin{notaprofesor}[Respuesta P9 — 4 pts]
  \begin{center}
  \begin{tabular}{lll}
    \toprule
    Paso & Expresión & Resultado \\
    \midrule
    1 & $30$ & \texttt{0001 1110} \\
    2 & $30 \;\texttt{>>}\; 3$ & \texttt{0000 0011} = 3 \\
    3 & \texttt{\textasciitilde{}(30 >> 3)} & \texttt{1111 1100} = $-4$ \\
    4 & $19$ & \texttt{0001 0011} \\
    5 & \texttt{(paso 3) \& 19} & \texttt{0001 0000} = \textbf{16} \\
    \bottomrule
  \end{tabular}
  \end{center}
  Resultado: \textbf{16}.

  \textbf{Calificación:} 1 pt por paso (desplazamiento, NOT, valor 19, AND final).
  El NOT es el más difícil; aceptar resultado correcto aunque el alumno omita
  mostrar todos los bits si escribe $\texttt{\textasciitilde 3} = -4 = \texttt{1111\,1100}$.
\end{notaprofesor}

\medskip
\textbf{P10.} Dado $n = 52$:
\begin{enumerate}
  \item ¿Cuánto vale \texttt{n \& (n-1)}?
  \item ¿Cuánto vale \texttt{n \& -n}?
  \item ¿Es $n$ una potencia de 2?
\end{enumerate}

\begin{notaprofesor}[Respuesta P10 — 3 pts]
  $52_{10} = \texttt{0011 0100}_2$.
  \begin{enumerate}
    \item \texttt{n - 1} = 51 = \texttt{0011 0011}.\\
          \texttt{0011 0100 \& 0011 0011 = 0011 0000} = $48$.\\
          Interpretación: se apagó el bit encendido de menor peso (bit 2, valor 4).
    \item $-52$ en 8 bits (complemento a 2): $256 - 52 = 204 = \texttt{1100 1100}$.\\
          \texttt{0011 0100 \& 1100 1100 = 0000 0100} = $4$.\\
          Confirmación: $4 = 2^2$, que es el bit encendido más bajo de 52.
    \item $n$ \textbf{no es potencia de 2} porque \texttt{n \& (n-1)} $\neq 0$
          (resultado fue 48, no 0).  Una potencia de 2 tiene exactamente un bit
          encendido; $52 = 32 + 16 + 4$ tiene tres.
  \end{enumerate}
\end{notaprofesor}

% -----------------------------------------------------------------------------
\subsection{Preguntas conceptuales}
% -----------------------------------------------------------------------------

\textbf{P11.} Explica con tus palabras por qué
\texttt{0.1 + 0.2 == 0.3} devuelve \texttt{false} en C++.

\begin{notaprofesor}[Respuesta P11 — 3 pts]
  Respuesta completa (3 pts):
  \begin{enumerate}
    \item $0.1_{10}$ es periódico en binario: $0.1_{10} = 0.0\overline{0011}_2$.
          No puede representarse exactamente con un número finito de bits.
    \item Al almacenarlo en un \texttt{float}/\texttt{double} (23 o 52 bits de mantisa),
          la cadena infinita se trunca, introduciendo un pequeño error de redondeo.
    \item $0.1 + 0.2$ produce la suma de \textit{dos} valores truncados.
          El resultado difiere de la representación truncada de $0.3_{10}$,
          que también es periódico y fue truncado de forma independiente.
  \end{enumerate}

  Respuesta aceptable (2 pts): menciona que el problema es que 0.1 no puede
  representarse exactamente en binario y que se acumula error de redondeo.

  Respuesta insuficiente (1 pt): solo dice ``error de redondeo'' sin explicar
  por qué existe ese error.
\end{notaprofesor}

\medskip
\textbf{P12.} ¿Por qué el rango del complemento a 2 no es simétrico?

\begin{notaprofesor}[Respuesta P12 — 2 pts]
  Con $n$ bits hay $2^n$ patrones posibles.
  El cero ocupa uno de los patrones del lado positivo
  (\texttt{0000\ldots0}).  Así, del total de $2^n$ patrones,
  $2^{n-1}$ tienen el bit de signo en 0 (cubre 0 a $+2^{n-1}-1$)
  y $2^{n-1}$ tienen el bit de signo en 1 (cubre $-2^{n-1}$ a $-1$).
  El lado negativo tiene un valor ``extra'' precisamente porque el cero
  no le ``roba'' ningún patrón.

  Consecuencia: con 8 bits existen $-128$ pero no $+128$;
  hay 128 valores negativos y solo 127 positivos (más el cero).
\end{notaprofesor}

\newpage

% =============================================================================
\section{Tiempos estimados por bloque}
% =============================================================================

\begin{center}
\begin{tabular}{lcc}
  \toprule
  \textbf{Bloque} & \textbf{Preguntas típicas} & \textbf{Tiempo estimado} \\
  \midrule
  Conversiones de bases       & 3–4 conversiones escalonadas & 18–22 min \\
  Complemento a 2             & 2–3 preguntas                & 12–15 min \\
  ASCII                       & 1–2 preguntas                &  8–10 min \\
  Fracciones e IEEE 754       & 2–3 preguntas                & 12–15 min \\
  Operadores de bits          & 2 expresiones paso a paso    & 10–12 min \\
  Técnicas y conceptual       & 2–3 preguntas                &  8–10 min \\
  \midrule
  \textbf{Total}              &                              & \textbf{68–84 min} \\
  \bottomrule
\end{tabular}
\end{center}

\begin{notaprofesor}[Gestión del tiempo]
  El examen está calibrado para terminarse en 80 minutos por un alumno
  que preparó bien el material.  Deja al menos 5 minutos de colchón.
  Si ves que el grupo se atrasa sistemáticamente en las conversiones,
  considera reducir el número de pasos a mostrar (e.g., solo pedir el resultado
  del bloque intermedio en hex$\to$bin, no todos los detalles).
\end{notaprofesor}

% =============================================================================
\section{Errores frecuentes y cómo abordarlos}
% =============================================================================

\begin{notaprofesor}[Los 5 errores más comunes]
  \begin{enumerate}
    \item \textbf{Complemento a 2: tratar el bit de signo como un peso positivo.}
          Un alumno ve \texttt{1000 0000} y calcula $128$ en lugar de $-128$.
          Remedio al calificar: si el proceso (invertir + sumar 1) está bien pero el signo
          final está equivocado, da crédito parcial.

    \item \textbf{Agrupación de bits desde la izquierda en lugar de la derecha.}
          En la conversión binario $\to$ hexadecimal, el grupo de la izquierda
          puede ser de menos de 4 bits si el número no es múltiplo de 4.
          Rellena con ceros por la izquierda, no por la derecha.

    \item \textbf{NOT bit a bit vs.\ NOT lógico.}
          Confunden \texttt{\textasciitilde{} 0} (que es $-1$, todos los bits a 1)
          con \texttt{!0} (que es \texttt{true} = 1).

    \item \textbf{Fracciones: copiar de memoria sin entender la regla.}
          Dicen ``$1/3$ no es exacta en binario'' sin poder justificarlo con
          los factores primos.  Si hay una fracción no estándar como $7/14$,
          fracasan.  Evalúa siempre con al menos una fracción que requiera
          simplificar primero ($7/14 = 1/2$ — exacta).

    \item \textbf{Precedencia entre \texttt{\&} y \texttt{|}.}
          Calculan \texttt{a \& b | c} como \texttt{a \& (b | c)} en lugar
          de \texttt{(a \& b) | c}.  Incluir al menos una pregunta donde esto
          haga diferencia numérica visible.
  \end{enumerate}
\end{notaprofesor}

% =============================================================================
\section{Criterios de calificación parcial}
% =============================================================================

\begin{notaprofesor}[Política general]
  \begin{itemize}
    \item \textbf{Proceso visible}: si el proceso es correcto pero hay un error
          aritmético puntual, otorga crédito parcial (típicamente 50--70\% de
          los puntos del ejercicio).
    \item \textbf{Solo respuesta}: si el alumno solo escribe el resultado sin
          mostrar proceso, no otorgas crédito aunque sea correcto (en preguntas
          que explícitamente pidan ``muestra el proceso'').  Anúncialo al inicio.
    \item \textbf{Respuestas conceptuales}: evalúa por ideas clave.  Una respuesta
          que menciona las ideas correctas aunque esté mal redactada vale más que
          una que copia la frase exacta de las notas pero no la comprende.
    \item \textbf{Verificación}: si el alumno hace una verificación y confirma
          un resultado correcto, suma puntos de proceso aunque los pasos
          intermedios sean menos detallados de lo esperado.
  \end{itemize}
\end{notaprofesor}

% =============================================================================
\section{Resumen del examen propuesto}
% =============================================================================

\begin{center}
\begin{tabular}{clcl}
  \toprule
  \textbf{\#} & \textbf{Pregunta} & \textbf{Pts} & \textbf{Tema} \\
  \midrule
   1 & Convierte 219 a binario                           & 2  & Conversiones \\
   2 & Convierte \texttt{6C3}$_{16}$ a octal             & 3  & Conversiones \\
   3 & Representa $-53$ en 8 bits complemento a 2        & 3  & Compl.\ a 2 \\
   4 & Interpreta \texttt{1010 0011} con signo            & 3  & Compl.\ a 2 \\
   5 & Decodifica \texttt{44 45 43 49 4D 41 4C} en ASCII  & 2  & ASCII \\
   6 & Mayúscula de \texttt{0x6D} con operaciones de bits & 3  & ASCII / bits \\
   7 & Convierte $\frac{1}{6}$ a binario                 & 3  & Fracciones \\
   8 & ¿Exactas en binario? $\frac{7}{16}$ y $\frac{3}{14}$ & 2 & Fracciones \\
   9 & Evalúa \texttt{(\textasciitilde(30>>3))\&19}       & 4  & Operadores \\
  10 & $n=52$: \texttt{n\&(n-1)}, \texttt{n\&-n}, ¿potencia de 2? & 3 & Técnicas \\
  11 & ¿Por qué \texttt{0.1+0.2 != 0.3}?                 & 3  & IEEE 754 \\
  12 & ¿Por qué el rango de compl.\ a 2 no es simétrico? & 2  & Conceptual \\
  \midrule
     & \textbf{Total}                                    & \textbf{33} & \\
  \bottomrule
\end{tabular}
\end{center}

\begin{notaprofesor}[Ajuste de puntuación]
  El total suma 33 puntos.  Puedes escalar a la calificación de tu
  institución (e.g., sobre 10: dividir entre 3.3).  También puedes
  agregar puntos extra con una pregunta de conversión adicional o un
  ejercicio de bit manipulation sobre los problemas de la sesión 07
  (e.g., verificar el estado de los enchufes para valores concretos
  de $n$ y $c$) para llegar a 40 puntos y tener mayor granularidad.
\end{notaprofesor}

\end{document}
